\documentclass[12pt,a4paper]{article}
\usepackage[utf8]{inputenc}
\usepackage[T1]{fontenc}
\usepackage{geometry}
\usepackage{enumitem}
\usepackage{fancyhdr}
\usepackage{lastpage}
\usepackage{hyperref}

\geometry{margin=2.5cm}
\pagestyle{fancy}
\fancyhf{}
\rhead{Page \thepage\ of \pageref{LastPage}}
\lhead{CSV402 - Quiz (en)}
\renewcommand{\headrulewidth}{0.4pt}

\title{Quiz: CSV402}
\author{Language: en}
\date{\today}

\begin{document}

\maketitle
\thispagestyle{fancy}

\section*{Instructions}
Answer all questions by selecting the best option (A, B, C, or D). The answer key is provided at the end of this document.

\bigskip


\subsection*{Question 1}
What is the main function of blockchain in the client-side validation paradigm used by RGB?

\begin{enumerate}[label=\Alph*., leftmargin=*]
\item Store all state transitions
\item Guarantee global data availability
\item Replicating smart contract data
\item Provide time stamping and prevent double spending
\end{enumerate}

\bigskip


\subsection*{Question 2}
What is an essential feature of the single use seals used in RGB?

\begin{enumerate}[label=\Alph*., leftmargin=*]
\item They guarantee the confidentiality of transactions
\item They can only be closed once
\item They can only be closed twice
\item They eliminate the need for consensus on blockchain
\end{enumerate}

\bigskip


\subsection*{Question 3}
In the CAP trilemma, what is the trade-off between customer-side validation and availability?

\begin{enumerate}[label=\Alph*., leftmargin=*]
\item State transitions are always published on Lightning
\item Each participant must keep the relevant data locally
\item Data is public and automatically replicated globally
\item Consistency sacrificed for availability
\end{enumerate}

\bigskip


\subsection*{Question 4}
What is the consequence of losing the stash in the RGB protocol?

\begin{enumerate}[label=\Alph*., leftmargin=*]
\item Contract history and associated data become irretrievable
\item Data is automatically replicated for all participants
\item A copy of the data is available on RGB's servers in Switzerland
\item Data can be reconstructed from the blockchain if you have the seed
\end{enumerate}

\bigskip


\subsection*{Question 5}
How does RGB improve scalability compared with traditional blockchain?

\begin{enumerate}[label=\Alph*., leftmargin=*]
\item By linking each contract directly to all UTXOs
\item By centralizing contract data on a few central servers
\item By limiting storage and validation to relevant state transitions only
\item By virtually increasing the size of Bitcoin blocks
\end{enumerate}

\bigskip


\subsection*{Question 6}
What method is used in RGB to define a single-use seal?

\begin{enumerate}[label=\Alph*., leftmargin=*]
\item A combination of several UTXOs
\item A specific outpoint
\item A specific Bitcoin address
\item A specific public key
\end{enumerate}

\bigskip


\subsection*{Question 7}
In RGB, where is the commitment when a single-use seal is closed?

\begin{enumerate}[label=\Alph*., leftmargin=*]
\item In the output of the transaction spending the outpoint
\item In the input of the transaction spending the outpoint
\item In the Witness of the transaction spending the outpoint
\item In the scriptSig of the transaction spending the outpoint
\end{enumerate}

\bigskip


\subsection*{Question 8}
How does RGB minimize its visibility?

\begin{enumerate}[label=\Alph*., leftmargin=*]
\item Key tweak
\item Sig tweak
\item Taproot (Tapret)
\item OP\_RETURN (Opret)
\end{enumerate}

\bigskip


\subsection*{Question 9}
What is the main purpose of deterministic Bitcoin commitments?

\begin{enumerate}[label=\Alph*., leftmargin=*]
\item Optimizing RGB multisignature generation
\item Reduce RGB transaction costs
\item Simplifying the structure of Bitcoin transactions
\item Ensure that a single commitment is inserted in a transaction
\end{enumerate}

\bigskip


\subsection*{Question 10}
On RGB, why is it recommended to use Tapret instead of Opret for commitments?

\begin{enumerate}[label=\Alph*., leftmargin=*]
\item Because an OP\_RETURN is not Lightning-compatible
\item To reduce transaction costs
\item To minimize visibility and avoid censorship
\item To maximize the size of commitments
\end{enumerate}

\bigskip


\subsection*{Question 11}
What is AluVM's main role in the RGB protocol?

\begin{enumerate}[label=\Alph*., leftmargin=*]
\item Analyze Bitcoin on-chain transactions in a virtual machine environment
\item Generate public keys for RGB users in a virtual machine environment
\item Execute scripts and contract operations in a virtual machine environment
\item Optimizing public blockchain validation in a virtual machine environment
\end{enumerate}

\bigskip


\subsection*{Question 12}
What makes up an Anchor in RGB?

\begin{enumerate}[label=\Alph*., leftmargin=*]
\item Genesis, Seal Definition, and Transition Bundle
\item Transaction ID, MPC, DBC, and Extra Transaction Proof
\item Signature, Transaction Hash, and Pedersen Commitment
\item Contract Schema, Business Logic, and Owned State
\end{enumerate}

\bigskip


\subsection*{Question 13}
What is an Assignment in the RGB protocol?

\begin{enumerate}[label=\Alph*., leftmargin=*]
\item A public update of the Global State
\item An output that modifies, updates or creates properties in the state of a contract
\item A unique cryptographic commitment stored in a Bitcoin transaction
\item A Bitcoin transaction partially signed before final validation
\end{enumerate}

\bigskip


\subsection*{Question 14}
What are the main operations of an RGB contract?

\begin{enumerate}[label=\Alph*., leftmargin=*]
\item State Transition, Genesis, and State Extension
\item Seal Definition, Commitment, and Validation
\item Genesis Block, UTXO Allocation, and Script Validation
\item Anchor Creation, DAG Update, and Transaction Signing
\end{enumerate}

\bigskip


\subsection*{Question 15}
What role does a Schema play in RGB?

\begin{enumerate}[label=\Alph*., leftmargin=*]
\item Validating Pedersen Commitments on Bitcoin transactions
\item Store consignments exchanged between participants
\item Create Seal Definitions for each UTXO
\item Define contract variables, rules and business logic
\end{enumerate}

\bigskip


\subsection*{Question 16}
What historical concept does RGB apply to the digital world?

\begin{enumerate}[label=\Alph*., leftmargin=*]
\item Peer-to-peer networks without validation
\item Bearer instruments
\item Centralized databases
\item Open registers accessible to all
\end{enumerate}

\bigskip


\subsection*{Question 17}
What are the two main components of a smart contract state in RGB?

\begin{enumerate}[label=\Alph*., leftmargin=*]
\item Global State and Owned States
\item Seal Definition and State Transition
\item Commitment and Genesis
\item Valency and Contract Rights
\end{enumerate}

\bigskip


\subsection*{Question 18}
What structure guarantees the temporal consistency of transitions in RGB?

\begin{enumerate}[label=\Alph*., leftmargin=*]
\item A simple linear chain
\item A dynamic peer-to-peer network
\item A binary tree for UTXOs
\item A DAG (Directed Acyclic Graph)
\end{enumerate}

\bigskip


\subsection*{Question 19}
What operations define the evolution of a smart contract in RGB?

\begin{enumerate}[label=\Alph*., leftmargin=*]
\item UTXO Transfer, Proof Generation, and Hash Commitment
\item Transition Bundles, Single-use Seals, and Valency Creation
\item Seal Creation, Contract Observance, and Final Validation
\item State Transition, Genesis, and State Extension
\end{enumerate}

\bigskip


\subsection*{Question 20}
What is the main purpose of smart contracts in the RGB environment?

\begin{enumerate}[label=\Alph*., leftmargin=*]
\item Simplifying the execution of ERC-20 tokens on Bitcoin
\item Implement automated Bitcoin deals without a centralized registry
\item Facilitating the creation of centralized Bitcoin registries
\item Provide a public interface for all Bitcoin transactions
\end{enumerate}

\bigskip


\subsection*{Question 21}
What is the key element that ensures Bob receives an asset during a State Transition in RGB?

\begin{enumerate}[label=\Alph*., leftmargin=*]
\item A control transaction without reference to a UTXO
\item A public statement from Alice
\item A new Seal Definition pointing to Bob's UTXO
\item A double signature shared by Alice and Bob
\end{enumerate}

\bigskip


\subsection*{Question 22}
What does a Consignment in RGB contain?

\begin{enumerate}[label=\Alph*., leftmargin=*]
\item A list of Bitcoin's public commitments
\item The customer-side data needed to validate a transition
\item A complete copy of the Bitcoin blockchain
\item Public keys of state transition participants
\end{enumerate}

\bigskip


\subsection*{Question 23}
Which RGB contract operation creates a new state without closing a seal?

\begin{enumerate}[label=\Alph*., leftmargin=*]
\item Valency Activation
\item Genesis
\item State Extension
\item State Transition
\end{enumerate}

\bigskip


\subsection*{Question 24}
How does RGB prevent double-spending during a State Transition?

\begin{enumerate}[label=\Alph*., leftmargin=*]
\item Validating Global State data
\item Closing an existing seal and creating a new one
\item By verifying all transactions on the blockchain
\item By imposing multi-sig signatures on every operation
\end{enumerate}

\bigskip


\subsection*{Question 25}
What mechanism in RGB allows multiple transitions to be grouped together in a single Bitcoin transaction?

\begin{enumerate}[label=\Alph*., leftmargin=*]
\item Seal Aggregation
\item Transition Bundle
\item State Consolidation
\item Multi-Commitment Tree
\end{enumerate}

\bigskip


\subsection*{Question 26}
What are the four components that define an RGB contract?

\begin{enumerate}[label=\Alph*., leftmargin=*]
\item Genesis, Interface, Interface Implementation and Extension
\item Genesis, Interface, Transition and Extension
\item Genesis, Schema, Transition and Interface Implementation
\item Genesis, Schema, Interface and Interface Implementation
\end{enumerate}

\bigskip


\subsection*{Question 27}
What is the purpose of Interface Implementation in RGB architecture?

\begin{enumerate}[label=\Alph*., leftmargin=*]
\item To store client-side cryptographic proofs
\item Compile AluVM scripts for validation
\item Map Schema fields to names defined in the Interface
\item To define the global states and owned states of the contract
\end{enumerate}

\bigskip


\subsection*{Question 28}
What is the object-oriented programming (OOP) equivalent of the Schema component in the RGB context?

\begin{enumerate}[label=\Alph*., leftmargin=*]
\item Class definition
\item Method prototyping
\item Implementing an interface
\item The class constructor
\end{enumerate}

\bigskip


\subsection*{Question 29}
In an RGB contract, what is the function of the Interface?

\begin{enumerate}[label=\Alph*., leftmargin=*]
\item Make contract data easy to read and manipulate
\item Store Genesis data on the customer's side
\item Running validation scripts on the blockchain
\item Define strict type system for reports
\end{enumerate}

\bigskip


\subsection*{Question 30}
What is the main role of Schema in an RGB contract?

\begin{enumerate}[label=\Alph*., leftmargin=*]
\item Describe the business logic and data structure of the contract
\item Defining the interface for RGB wallets
\item Managing the compilation of scripts on the blockchain
\item On-chain cryptographic validation
\end{enumerate}

\bigskip


\subsection*{Question 31}
What are the components contained in an RGB consignment that Bob (recipient) receives for a contract?

\begin{enumerate}[label=\Alph*., leftmargin=*]
\item Genesis, Interface, PSBT and History
\item Genesis, Schema, Transition and Extension
\item Schema, Interface, Signature and PSBT
\item Genesis, Schema, Interface and Interface Implementation
\end{enumerate}

\bigskip


\subsection*{Question 32}
What can I do with an invoice in an RGB transfer?

\begin{enumerate}[label=\Alph*., leftmargin=*]
\item Indicate the type of transaction, the quantity of assets required and the recipient's seal definition
\item Generate a new Schema to modify the asset's business logic
\item Building a PSBT to replace the Bitcoin transaction
\item Create a new Genesis for the asset
\end{enumerate}

\bigskip


\subsection*{Question 33}
What should Alice (sender) do after receiving Bob's invoice?

\begin{enumerate}[label=\Alph*., leftmargin=*]
\item Build a PSBT with its UTXO and generate the transfer consignment terminal
\item Import the consignment into your own wallet for contract details
\item Issue a new invoice for an asset issue
\item Sign the Bitcoin transaction that spends its UTXO containing the single-use seal
\end{enumerate}

\bigskip


\subsection*{Question 34}
What is the purpose of terminal consignment in the RGB transfer process?

\begin{enumerate}[label=\Alph*., leftmargin=*]
\item Delete old transitions
\item Provide State Transitions history
\item Validate Genesis contract
\item Generate a new Schema for the transfer
\end{enumerate}

\bigskip


\subsection*{Question 35}
What action does Bob (receiver) perform once he receives the terminal consignment from Alice (payer)?

\begin{enumerate}[label=\Alph*., leftmargin=*]
\item Create an invoice for a new operation
\item Sign Alice's PSBT and distribute it onchain
\item Check validity of State Transition and sign
\item Modify the Contract Schema and send it to Alice
\end{enumerate}

\bigskip


\subsection*{Question 36}
How do I install the rgb tool via Cargo?

\begin{enumerate}[label=\Alph*., leftmargin=*]
\item cargo install rgb --features all
\item cargo install rgb-contracts --all-features
\item cargo install rgb-bitcoin --all
\item cargo build rgb-lightning --release
\end{enumerate}

\bigskip


\subsection*{Question 37}
What is the role of the stash?

\begin{enumerate}[label=\Alph*., leftmargin=*]
\item Keep all client-side data locally
\item Compile and store all generated invoices
\item Managing Bitcoin private keys
\item Automatically publish transactions on the blockchain
\end{enumerate}

\bigskip


\subsection*{Question 38}
How do I import the NonInflatableAsset.rgb Schema into the local stash?

\begin{enumerate}[label=\Alph*., leftmargin=*]
\item rgb clone schemata/NonInflatableAssets.rgb
\item rgb import schemata/NonInflatableAssets.rgb
\item rgb import schema/NonInflatableAssets.rgb
\item rgb clone schema/NonInflatableAssets.rgb
\end{enumerate}

\bigskip


\subsection*{Question 39}
In the YAML file used to issue an asset, which field defines the total number of tokens issued?

\begin{enumerate}[label=\Alph*., leftmargin=*]
\item assetOwner
\item amount
\item issuedSupply
\item spec
\end{enumerate}

\bigskip


\subsection*{Question 40}
Which command exports a contract from the stash to an external file?

\begin{enumerate}[label=\Alph*., leftmargin=*]
\item rgb dump '<ContractId>' myContractPBN.rgb
\item rgb download '<ContractId>' myContractPBN.rgb
\item rgb export '<ContractId>' myContractPBN.rgb
\item rgb issue '<ContractId>' myContractPBN.rgb
\end{enumerate}

\bigskip


\subsection*{Question 41}
What is the main aim of integrating RGB into the Lightning Network?

\begin{enumerate}[label=\Alph*., leftmargin=*]
\item Executing RGB contracts directly on the Bitcoin blockchain
\item Replace conventional Lightning transactions with RGB smart contracts
\item Transfer RGB assets via off-chain channels
\item Increase channel capacity by adding new UTXOs
\end{enumerate}

\bigskip


\subsection*{Question 42}
What elements are needed to create a Lightning channel carrying RGB assets?

\begin{enumerate}[label=\Alph*., leftmargin=*]
\item Funding in bitcoins and funding in RGB assets
\item RGB-only funding to start with
\item Bitcoin-only funding to start with
\item No funding required
\end{enumerate}

\bigskip


\subsection*{Question 43}
In the event of payment within the RGB channel, what happens to the RGB outputs in the commitment transactions?

\begin{enumerate}[label=\Alph*., leftmargin=*]
\item They are removed and replaced by new multi-sig
\item They are updated to reflect the redistribution of assets among participants
\item They are converted into classic Bitcoin UTXOs
\item They remain unchanged to preserve the initial history until the channel is closed
\end{enumerate}

\bigskip


\subsection*{Question 44}
What does the "RGB state transition" represent in the context of the Lightning Network?

\begin{enumerate}[label=\Alph*., leftmargin=*]
\item On-chain validation script replaced by HTLC
\item A new state of multisig 2/2 onchain
\item Updated RGB asset allocation integrated into commitment transactions
\item Updated RGB asset allocation integrated into the funding transaction
\end{enumerate}

\bigskip


\subsection*{Question 45}
How is the safety of an RGB transition in a Lightning channel ensured if an obsolete state is broadcast?

\begin{enumerate}[label=\Alph*., leftmargin=*]
\item Through the intervention of the RGB asset issuer via its signature associated with the Genesis
\item Through the punitive mechanics of the RGB protocol
\item Through the punitive mechanics of the Lightning Channel
\item By automatically cancelling all previous transactions
\end{enumerate}

\bigskip


\subsection*{Question 46}
Which RGB standard is used to define fungible tokens?

\begin{enumerate}[label=\Alph*., leftmargin=*]
\item RGB21
\item RGB22
\item RGB25
\item RGB20
\end{enumerate}

\bigskip


\subsection*{Question 47}
Which RGB standard allows you to create unique assets or "NFTs"?

\begin{enumerate}[label=\Alph*., leftmargin=*]
\item RGB25
\item RGB30
\item RGB21
\item RGB20
\end{enumerate}

\bigskip


\subsection*{Question 48}
What technology can be used to merge different versions of the stash in the event of a discrepancy?

\begin{enumerate}[label=\Alph*., leftmargin=*]
\item Log-Based Replication
\item HTLC
\item CRDT
\item Single-use seal
\end{enumerate}

\bigskip


\subsection*{Question 49}
What is the fundamental concept behind the RGB protocol?

\begin{enumerate}[label=\Alph*., leftmargin=*]
\item Customer-side validation via on-chain cryptographic commitments
\item On-chain validation by Turing-complete smart contracts
\item Exclusive use of unencrypted Bitcoin transactions
\item Contract execution on a blockchain dedicated to smart contracts
\end{enumerate}

\bigskip


\subsection*{Question 50}
How does RGB ensure confidentiality in digital asset management?

\begin{enumerate}[label=\Alph*., leftmargin=*]
\item By hiding users' identities via decentralized mixers
\item By registering only cryptographic commitments on the blockchain
\item Using only temporary Bitcoin addresses for each transfer
\item By encrypting all end-to-end transactions on the blockchain
\end{enumerate}

\bigskip


\subsection*{Question 51}
How are cryptographic commitments (via Taproot or OP\_RETURN) used in RGB?

\begin{enumerate}[label=\Alph*., leftmargin=*]
\item They anchor state transitions in the blockchain
\item They store contract data on the blockchain
\item They create nodes for transaction verification
\item They replace client-side validation with on-chain smart contract execution
\end{enumerate}

\bigskip


\subsection*{Question 52}
What is a "Seal Definition" in the context of RGB Assignment?

\begin{enumerate}[label=\Alph*., leftmargin=*]
\item The grouping of data exchanged between the parties, subject to Client-side Validation
\item Compiling contracts in off-chain executable code
\item The DAG branch of the state transition history
\item The reference to a UTXO that establishes ownership of an Owned State
\end{enumerate}

\bigskip


\subsection*{Question 53}
What is client-side validation in the RGB protocol?

\begin{enumerate}[label=\Alph*., leftmargin=*]
\item Local verification of contract data and transitions
\item Automatic reset of contract histories in the event of discrepancies
\item Exclusive use of centralized nodes to validate transactions
\item Executing smart contracts directly on the Bitcoin blockchain
\end{enumerate}

\bigskip


\subsection*{Question 54}
How does RGB pave the way for DEX layer-2?

\begin{enumerate}[label=\Alph*., leftmargin=*]
\item By enabling inter-blockchain asset exchange via atomic swaps
\item By enabling the use of a centralized exchange for all RGB transactions
\item By enabling the exchange of off-chain assets via Lightning channels and atomic swaps
\item By enabling the exchange of assets via on-chain smart contracts
\end{enumerate}

\bigskip


\subsection*{Question 55}
What is RGBlib in the RGB ecosystem?

\begin{enumerate}[label=\Alph*., leftmargin=*]
\item An Android wallet dedicated to RGB asset management
\item An alternative client for the RGB protocol
\item A website to display RGB NFTs
\item A library for creating applications managing RGB assets
\end{enumerate}

\bigskip


\subsection*{Question 56}
What is rgb-lightning-node in the RGB ecosystem?

\begin{enumerate}[label=\Alph*., leftmargin=*]
\item A mobile wallet to manage RGB assets
\item A daemon for managing RGB-compatible Lightning channels
\item A mining tool for RGB assets on Bitcoin
\item A proxy server to exchange RGB logs
\end{enumerate}

\bigskip


\subsection*{Question 57}
What are the main components required to run an RGB-compatible Lightning node?

\begin{enumerate}[label=\Alph*., leftmargin=*]
\item An RGB mobile wallet, an RGB proxy server and a bitcoind
\item A bitcoind, an indexer and an RGB proxy server
\item A Lightning node, an indexer and an IPFS server
\item Web client, indexer and RGB proxy server
\end{enumerate}

\bigskip


\subsection*{Question 58}
On RLN, which command installs the "rgb-lightning-node" binary after cloning the repository?

\begin{enumerate}[label=\Alph*., leftmargin=*]
\item cargo build --release --path .
\item cargo install --locked --debug --path .
\item cargo run --locked --path .
\item cargo deploy --debug --path .
\end{enumerate}

\bigskip


\subsection*{Question 59}
What can curl -X POST /issueassetnia do on RLN?

\begin{enumerate}[label=\Alph*., leftmargin=*]
\item Create a new fungible inflationary RGB asset
\item Create a new semi-fungible RGB asset
\item Create a new non-fungible RGB asset
\item Create a new fungible, non-inflationary RGB asset
\end{enumerate}

\bigskip


\subsection*{Question 60}
What is the function of the RGB proxy server in an RGB Lightning Node environment?

\begin{enumerate}[label=\Alph*., leftmargin=*]
\item Avoid duplicate tickers for RGB tokens
\item Provide public keys for nodes to open channels
\item Keep history of RGB state transitions
\item Facilitating the exchange of RGB consignments
\end{enumerate}

\bigskip


\newpage
\section*{Answer Key}

\begin{enumerate}
\item \textbf{Question 1:} D
\\ \textit{The Bitcoin blockchain is used in RGB to time-stamp commitments and prevent double-spending, without storing all state transitions.}

\item \textbf{Question 2:} B
\\ \textit{Single use seals enable the creation of unique cryptographic commitments, preventing duplication or double spending.}

\item \textbf{Question 3:} B
\\ \textit{In client-side validation, global availability is not guaranteed, as participants have to manage their own data.}

\item \textbf{Question 4:} A
\\ \textit{In client-side validation, each user is responsible for his or her stash, which is not replicated elsewhere, making its loss definitive.}

\item \textbf{Question 5:} C
\\ \textit{RGB reduces storage and validation requirements by retaining only the data relevant to each participant.}

\item \textbf{Question 6:} B
\\ \textit{RGB defines a single-use seal as a specific outpoint, which can then be closed by the expenditure of this outpoint (TXID + output number).}

\item \textbf{Question 7:} A
\\ \textit{The commitment is inserted into the output of the transaction that spends the outpoint defining the single-use seal.}

\item \textbf{Question 8:} C
\\ \textit{Taproot (Tapret) enables commitments to be inserted into a transaction discreetly and efficiently, while ensuring compatibility with current standards.}

\item \textbf{Question 9:} D
\\ \textit{Bitcoin deterministic commitments guarantee that a single commitment is inserted into a transaction, to avoid hidden competing commitments.}

\item \textbf{Question 10:} C
\\ \textit{Tapret lets you integrate commitments discreetly, unlike OP\_RETURN, which makes them immediately visible and subject to censorship.}

\item \textbf{Question 11:} C
\\ \textit{AluVM is designed for smart contract validation and distributed computing within the RGB framework.}

\item \textbf{Question 12:} B
\\ \textit{An Anchor contains the data needed to prove the inclusion of a unique commitment in a Bitcoin transaction.}

\item \textbf{Question 13:} B
\\ \textit{An Assignment combines a Seal Definition and an Owned State to represent an allocated portion of the state.}

\item \textbf{Question 14:} A
\\ \textit{These operations modify or complete the status of a contract according to its Schema.}

\item \textbf{Question 15:} D
\\ \textit{The Schema structures the state of the contract and determines the permitted transitions and validation conditions.}

\item \textbf{Question 16:} B
\\ \textit{RGB transposes the idea of bearer instruments to digital.}

\item \textbf{Question 17:} A
\\ \textit{The Global State contains public data, while Owned States are associated with specific holders via UTXOs.}

\item \textbf{Question 18:} D
\\ \textit{A DAG enables state transitions to be structured, while allowing parallel evolutions without inconsistencies.}

\item \textbf{Question 19:} D
\\ \textit{These operations create or modify the state of the contract, while guaranteeing its validity via Bitcoin anchors.}

\item \textbf{Question 20:} B
\\ \textit{RGB smart contracts enable the automation of agreements while avoiding the use of centralized registers, relying on Bitcoin for cryptographic anchors.}

\item \textbf{Question 21:} C
\\ \textit{The Seal Definition cryptographically links the asset to Bob's UTXO, ensuring that he owns it.}

\item \textbf{Question 22:} B
\\ \textit{The Consignment includes the RGB information enabling the receiver to validate the state transition according to the contract rules.}

\item \textbf{Question 23:} C
\\ \textit{The State Extension allows you to create a new state or rights (Valencies), without the seal being immediately closed.}

\item \textbf{Question 24:} B
\\ \textit{Closing a seal with a spent UTXO proves that a transition is unique and irrevocable.}

\item \textbf{Question 25:} B
\\ \textit{A Transition Bundle is a set of RGB State Transitions (belonging to the same contract) which are all involved in the same Bitcoin witness transaction.}

\item \textbf{Question 26:} D
\\ \textit{An RGB contract is made up of four distinct elements; Genesis, Schema, Interface and Interface Implementation.}

\item \textbf{Question 27:} C
\\ \textit{Interface Implementation establishes the link between the Schema's business logic and the presentation readable by the Interface.}

\item \textbf{Question 28:} A
\\ \textit{The Schema corresponds to the definition of a class in OOP, which describes the properties, methods and rules of the business logic.}

\item \textbf{Question 29:} A
\\ \textit{The aim of the Interface is to present contract information in a way that is clear to users and software, without affecting the business logic.}

\item \textbf{Question 30:} A
\\ \textit{The Schema defines the declarative structure, data types and validation rules of the contract, enabling you to manage the business logic on the customer's side.}

\item \textbf{Question 31:} D
\\ \textit{The consignment of an RGB contract includes the Genesis, Schema, Interface and Interface Implementation, which provide all the data needed to validate the contract.}

\item \textbf{Question 32:} A
\\ \textit{The invoice is used to inform the issuer of the transaction of the type of operation, the quantity of assets to be transferred and the destination (seal definition).}

\item \textbf{Question 33:} A
\\ \textit{Alice must create a PSBT to spend the UTXO associated with the asset, then generate a terminal consignment that includes the new transition to Bob.}

\item \textbf{Question 34:} B
\\ \textit{The terminal consignment collects the history of the contract to prove the legitimacy of the asset transferred to date.}

\item \textbf{Question 35:} C
\\ \textit{Bob examines the consignment, makes sure that the transition complies with the rules, and signs off to confirm his agreement.}

\item \textbf{Question 36:} B
\\ \textit{Installation is performed with the command "cargo install rgb-contracts --all-features", the crate being named "rgb-contracts", which installs the rgb tool.}

\item \textbf{Question 37:} A
\\ \textit{The stash is the set of client-side data that a user stores locally for one or more RGB contracts, in order to perform client-side validation.}

\item \textbf{Question 38:} B
\\ \textit{The compiled Schema is imported into the local stash using the "rgb import schemata/NonInflatableAssets.rgb" command.}

\item \textbf{Question 39:} C
\\ \textit{The "issuedSupply" field in the YAML file defines the total amount of tokens issued for the contract.}

\item \textbf{Question 40:} C
\\ \textit{The "rgb export" command is used to extract a contract from the stash and save it as a file for easy sharing.}

\item \textbf{Question 41:} C
\\ \textit{RGB integrated into the Lightning Network allows RGB assets to be moved via state transitions embedded in commitment transactions, while remaining off-chain.}

\item \textbf{Question 42:} A
\\ \textit{The channel requires both bitcoin funding on a 2/2 multisig, as well as RGB asset funding to allocate assets on the same channel.}

\item \textbf{Question 43:} B
\\ \textit{When a payment is made, the RGB state transition adjusts the RGB outputs to redistribute the assets according to the new channel balance.}

\item \textbf{Question 44:} C
\\ \textit{The RGB state transition updates the distribution of RGB assets in the channel using commitment transactions.}

\item \textbf{Question 45:} C
\\ \textit{The Lightning Network's classic punitive mechanism ensures that the dissemination of an old condition can be punished by the recovery of all funds involved, including RGB assets.}

\item \textbf{Question 46:} D
\\ \textit{RGB20 is the standard used to create and manage fungible tokens over the RGB protocol.}

\item \textbf{Question 47:} C
\\ \textit{The RGB21 standard is designed to manage unique assets, enabling private metadata to be associated with NFTs.}

\item \textbf{Question 48:} C
\\ \textit{CRDTs ensure data consistency by enabling multiple versions of the stash to be merged without client-side conflicts.}

\item \textbf{Question 49:} A
\\ \textit{RGB is based on a client-side validation model where commitments (via Taproot or OP\_RETURN) are entered on-chain, while most of the logic and data is kept off-chain.}

\item \textbf{Question 50:} B
\\ \textit{By storing only on-chain cryptographic evidence, RGB keeps sensitive contract and transfer information confidential.}

\item \textbf{Question 51:} A
\\ \textit{Commitments enable status updates to be securely linked to a UTXO, in order to prove the existence of a certain piece of information at a given time, without revealing the information itself.}

\item \textbf{Question 52:} D
\\ \textit{The "Seal Definition" links a UTXO to a specific part of the contract state, guaranteeing the correct allocation of a right or asset.}

\item \textbf{Question 53:} A
\\ \textit{Client-side validation enables each participant to independently verify that the data received complies with the rules defined by the RGB contract.}

\item \textbf{Question 54:} C
\\ \textit{The ability to carry out atomic swaps and exchanges via Lightning channels offers the possibility of creating decentralized DEXs, without the disadvantages of on-chain exchanges.}

\item \textbf{Question 55:} D
\\ \textit{RGBlib offers a simplified interface for developers, enabling them to manage the emission, transfer and validation of RGB assets without having to master all the technical details.}

\item \textbf{Question 56:} B
\\ \textit{rgb-lightning-node is a Rust daemon designed to integrate RGB assets into Lightning channels by adding RGB state transitions to engagement transactions.}

\item \textbf{Question 57:} B
\\ \textit{To run an RGB-enabled Lightning node, you need a bitcoind node to interact with the Bitcoin blockchain, an indexer to mine on-chain transactions, and an RGB proxy server to facilitate the exchange of consignments between peers.}

\item \textbf{Question 58:} B
\\ \textit{The "cargo install --locked --debug --path ." command compiles and installs the "rgb-lightning-node" binary.}

\item \textbf{Question 59:} D
\\ \textit{The "/issueassetnia" command is used to issue a new RGB NIA ("Non Inflatable Asset") token.}

\item \textbf{Question 60:} D
\\ \textit{The RGB proxy server simplifies the communication of RGB consignments between Lightning participants by acting as an intermediary, enabling a smoother user experience close to that currently experienced with Bitcoin and Lightning.}


\end{enumerate}

\end{document}
